% Drawing 138/06-23 -- folder 138, batch 06, archive page 114 (PDF page 115).
% Transcription: un tripode ; une boucle avec une queue ; un segment à deux sommets prolongé.
% Kind: small maps (case 3). Method 2 (schematic redraw from the facsimile). Status: ok.
% Notes: Case 3: a tripod; a loop with a tail; a segment with two vertices, prolonged (brace under the last two). The symmetry groups and the types written under the maps are text of the transcription and are not drawn. Small planar maps listed by degree (number of edges) in the text; redrawn, not encoded (each would be a map/1 with at most a few edges; the text gives the classification).
% Coordinates: 1 unit = 1 mm = 1/100 of the crop width (bbox in page fractions: [0.189, 0.173, 0.567, 0.25]).
% Redrawn by Claude (claude-opus-5-5), issue #73, 2026-11-07; not checked by a human; not compiled (no TeX in the session).
% Rights: reproduces part of the work -- behind RENDER_TRANSCRIPTS, never published by itself.
\documentclass[tikz,border=3pt]{standalone}
\usepackage{amsmath,amssymb}
\usetikzlibrary{patterns,arrows.meta}
\tikzset{pencil/.style={black!45}, struck/.style={black!35}, mark/.style={circle,draw,inner sep=1.3pt}, pt/.style={circle,fill,inner sep=0.9pt}, opt/.style={circle,draw,fill=white,inner sep=1.1pt}}
\begin{document}
\begin{tikzpicture}[x=1mm,y=1mm,line cap=round,line join=round,font=\scriptsize]
\draw (10,-10) -- (10,-2.2);
\draw (10,-10) -- (4.4,-15.6);
\draw (10,-10) -- (16.7,-14.4);
\node[pt] at (10,-10) {};
\draw (46.7,-10) circle[radius=7.8];
\draw (40,-12.2) -- (31.1,-12.2);
\node[pt] at (40,-12.2) {};
\draw (72.2,-12.2) -- (96.7,-12.2);
\node[pt] at (72.2,-12.2) {};
\node[pt] at (84.4,-12.2) {};
\draw plot[smooth, tension=0.55] coordinates {(48.9,-22.9) (51.1,-25.6) (61.1,-25.6) (63.3,-28.2) (65.6,-25.6) (75.6,-25.6) (77.8,-22.9)};
\end{tikzpicture}
\end{document}
